In this project, the use of rich-burn, quick-mix, lean-burn (RQL) combustion to reduce \ce{NOx} emissions in amonia-fueled systems. Generally, ammonia is a promising clean fuel. However, it poses challenges due to high \ce{NOx} formation in lean premixed combustion. To address this issue, RQL provides a solution by initiating combustion in a fuel-rich environment that suppresses \ce{NOx} formation with a rapid air mixing and a final combustion in a lean zone. This combustion process is modeled in Cantera by coupling 3 reactors in series: a combustor rich burn zone, a mixer with secondary air addition, a combustor lean burn zone. In addition to this, numerical simulation is done using Ansys that uses a simplified RQL combustion geometry with defined zones for fuel injection, air mixing, and lean burn. Performance is then evaluated on the basis of emissions. This study shows the potential of RQL combustion that reduces \ce{NOx} emissions.